\chapter{Male health}

\section*{Definition and Overview}
Male health refers to a broad spectrum of medical, physical, and psychological issues that uniquely or predominantly affect men and boys throughout their lifespan. It encompasses male reproductive health (such as prostate and testicular conditions), sexual dysfunction, hormonal changes, cardiovascular wellness, and mental health. Globally, male health is recognised as a critical focus area due to distinct biological differences and behavioural patterns that lead to shorter life expectancies, higher rates of chronic diseases, and a lower propensity to seek medical care compared to women. Comprehensive male health management emphasises early screening, lifestyle modification, and addressing barriers to healthcare access.

\section*{Epidemiology}
Men experience unique epidemiological patterns that vary by age, socioeconomic status, and geography. In Australia, the average life expectancy for males is approximately 81.3 years, which is roughly 4 years shorter than that of females. The leading causes of death among men are ischaemic heart disease, lung cancer, dementia, and cerebrovascular disease. Prostate cancer is the most commonly diagnosed non-skin cancer in Australian men, with over 24,000 new cases diagnosed annually. Suicide is a major concern, with men accounting for approximately 75\% of all suicide deaths in Australia, particularly among young and middle-aged cohorts. Chronic conditions such as obesity (affecting over 70\% of Australian men) and type 2 diabetes are highly prevalent and drive cardiovascular morbidity.

\section*{Aetiology and Pathophysiology}
The determinants of male health conditions are multifactorial, spanning biological, genetic, lifestyle, and psychosocial factors:
\begin{itemize}
    \item \textbf{Genetic and hormonal factors:} Inherent susceptibility to X-linked disorders, male pattern baldness, and the age-related decline in testosterone levels (hypogonadism) that impacts muscle mass, bone density, and metabolic pathways.
    \item \textbf{Cardiovascular biology:} Elevated risks of atherosclerosis due to protective oestrogen deficit compared to pre-menopausal women, combined with vascular endothelial dysfunction.
    \item \textbf{Oncological changes:} Age-dependent mutations leading to benign prostatic hyperplasia (BPH) or adenocarcinoma of the prostate, and germ cell mutations in younger men leading to testicular cancer.
    \item \textbf{Psychosocial stressors:} Societal expectations of masculinity ("masculine norm adherence") that discourage vulnerability, leading to underreporting of mental health symptoms and delayed healthcare-seeking behaviours.
    \item \textbf{Lifestyle risk factors:} High prevalence of smoking, hazardous alcohol consumption, poor dietary choices, and physical inactivity that accelerate metabolic syndrome and visceral adiposity.
\end{itemize}

\section*{Clinical Features}
Common presentations across different domains of male health include:
\begin{itemize}
    \item \textbf{Urinary symptoms:} Frequency, nocturia, hesitancy, weak stream, or incomplete emptying (often due to BPH or prostate enlargement).
    \item \textbf{Sexual and reproductive dysfunction:} Erectile dysfunction, reduced libido, premature ejaculation, or infertility.
    \item \textbf{Testicular changes:} Painless lumps, swelling, heaviness, or acute pain in the scrotum.
    \item \textbf{Cardiometabolic features:} Central obesity (increased waist circumference), elevated blood pressure, fatigue, and chest pain upon exertion.
    \item \textbf{Mental health symptoms:} Irritability, anger, withdrawal, risk-taking behaviour, substance abuse, and physical symptoms like headaches or sleep issues (often masking underlying depression).
\end{itemize}

\section*{Differential Diagnosis}
Clinical complaints must be carefully investigated to identify the exact cause:
\begin{itemize}
    \item \textbf{Prostate cancer vs. BPH:} Both cause urinary symptoms and can elevate PSA, but require different diagnostic and treatment paradigms.
    \item \textbf{Depression vs. Low testosterone:} Fatigue, low mood, and reduced libido can be caused by either condition or both concurrently.
    \item \textbf{Erectile dysfunction etiologies:} Distinguish between psychogenic causes (performance anxiety, stress) and organic causes (atherosclerosis, diabetes, neuropathy).
    \item \textbf{Testicular mass vs. Epididymitis vs. Hydrocele:} Differentiate neoplastic lesions from inflammatory processes or benign fluid collections.
    \item \textbf{Angina vs. Gastroesophageal reflux:} Distinguish life-threatening myocardial ischaemia from benign acid reflux.
\end{itemize}

\section*{When to Seek Medical Care}
Men should undergo regular health check-ups and seek medical attention for new, progressive, or concerning symptoms.

\textbf{Call 000 (or local emergency services) for:}
\begin{itemize}
    \item \textbf{Chest pain or pressure:} Sudden pain spreading to the arm, neck, or jaw, accompanied by sweating or shortness of breath (suspected myocardial infarction).
    \item \textbf{Acute scrotal pain:} Sudden, severe testicular pain, which may indicate testicular torsion requiring emergency surgery within hours.
    \item \textbf{Severe mental health crisis:} Direct expressions of suicidal intent or self-harm.
    \item \textbf{Sudden neurological deficits:} Face drooping, arm weakness, or speech difficulty (suspected stroke).
\end{itemize}
Other reasons to see a doctor include a painless lump in the testicles, persistent changes in bowel or urinary habits, or chronic feelings of hopelessness.

\section*{Diagnosis and Investigations}
Screening and diagnostic tests play a vital role in managing male health:
\begin{itemize}
    \item \textbf{Prostate-Specific Antigen (PSA) test:} A blood test used to screen for and monitor prostate cancer, in combination with a digital rectal examination (DRE).
    \item \textbf{Cardiovascular risk assessment:} Regular checks of blood pressure, fasting lipid profiles (cholesterol), and blood glucose (HbA1c).
    \item \textbf{Scrotal ultrasound:} The primary imaging modality used to evaluate testicular lumps or scrotal pain.
    \item \textbf{Hormonal profile:} Early morning blood tests for total testosterone, SHBG, LH, and FSH to investigate suspected hypogonadism or infertility.
    \item \textbf{Mental health screening tools:} Questionnaires like the Patient Health Questionnaire (PHQ-9) adapted to identify depression in men.
    \item \textbf{Bowel cancer screening:} Immunochemical faecal occult blood tests (iFOBT) recommended every two years for men aged 50 to 74.
\end{itemize}

\section*{Management}
Intervention requires a multidimensional approach addressing physical, lifestyle, and psychological domains:
\begin{itemize}
    \item \textbf{Medical and pharmacological management:} Pharmacotherapy for BPH (5-alpha-reductase inhibitors or alpha-blockers), erectile dysfunction (PDE5 inhibitors), hypertension, dyslipidaemia, and testosterone replacement therapy (TRT) for hypogonadism.
    \item \textbf{Surgical options:} Transurethral resection of the prostate (TURP) for severe BPH, orchidectomy for testicular cancer, or coronary revascularisation for heart disease.
    \item \textbf{Psychosocial support:} Counselling, cognitive behavioural therapy (CBT), and community-based peer support groups (e.g. Men's Sheds) to improve mental health and reduce isolation.
    \item \textbf{Lifestyle interventions:} Weight management, structured physical activity (at least 150 minutes of moderate exercise per week), smoking cessation programmes, and reducing alcohol intake.
\end{itemize}

\section*{Complications}
Untreated male health issues can lead to significant complications:
\begin{itemize}
    \item \textbf{Cardiovascular events:} Myocardial infarction, stroke, or congestive heart failure secondary to unmanaged metabolic syndrome.
    \item \textbf{Advanced oncological disease:} Metastatic prostate or testicular cancer due to delayed presentation.
    \item \textbf{Urinary retention and renal failure:} Chronic bladder outlet obstruction from untreated BPH.
    \item \textbf{Chronic depression and suicide:} Unaddressed mental health struggles leading to severe crisis.
    \item \textbf{Osteoporotic fractures:} Bone loss secondary to untreated chronic hypogonadism.
\end{itemize}

\section*{Prognosis and Outcomes}
The prognosis for most male health conditions is highly favourable when detected and treated early. Diseases such as testicular cancer have a 5-year survival rate exceeding 95\% if diagnosed at an early stage. Active management of cardiovascular risk factors through lifestyle and medication can add years of high-quality life and prevent premature death. Encouraging open communication, reducing societal stigma around mental health, and promoting regular screening are the most effective means to improve long-term outcomes for men.

\section*{Prevention and Self-Care}
Prevention and self-care are essential to maintaining long-term health and well-being:
\begin{itemize}
    \item \textbf{Perform testicular self-examinations:} Monthly checks of the testicles for lumps, changes in size, or tenderness.
    \item \textbf{Adopt a heart-healthy diet:} Focus on a Mediterranean-style diet high in vegetables, fruits, whole grains, and healthy fats while limiting red meat and processed foods.
    \item \textbf{Stay active:} Aim for a combination of aerobic and strength training exercises throughout the week.
    \item \textbf{Get screened regularly:} Visit a general practitioner for scheduled blood pressure, blood glucose, cholesterol, and bowel cancer checks.
    \item \textbf{Connect with others:} Maintain strong social relationships and speak to trusted friends or professionals when feeling overwhelmed.
\end{itemize}

\section*{Sources and Further Reading}
The following reputable organisations provide further information on male health:
\begin{itemize}
    \item Healthy Male (Andrology Australia). \textit{Information on reproductive and sexual health for men.} https://www.healthymale.org.au/
    \item Beyond Blue. \textit{Men's mental health and depression support.} https://www.beyondblue.org.au/
    \item Better Health Channel. \textit{Men's health screening checklist and guidelines.} https://www.betterhealth.vic.gov.au/
    \item Mayo Clinic. \textit{Men's health overview, disease prevention, and healthy living.} https://www.mayoclinic.org/
    \item Healthdirect Australia. \textit{Male health topics and clinical advice.} https://www.healthdirect.gov.au/
\end{itemize}
